\selectlanguage{ngerman}%
\frenchspacing%
\begin{ledgroupsized}[r]{120mm}%
\footnotesize%
\pstart%
\noindent%
\textbf{Überlieferung:}%
\pend%
\end{ledgroupsized}%
\begin{ledgroupsized}[r]{114mm}%
\footnotesize%
\pstart%
\parindent -6mm%
\makebox[6mm][l]{\textit{L}}%
Aufzeichnung:
LH XXXVII\;4 Bl.\;93.
Ein Blatt 16\textsuperscript{o},
Ränder beschnitten und zum Teil gerissen (mit geringfügigem Textverlust);
Papiererhaltungsmaßnahmen.
Zwei vollbeschriebene Seiten.
\pend%
\end{ledgroupsized}%
\vspace{5mm}%
\begin{ledgroup}%
\footnotesize%
\pstart%
\noindent%
\textbf{Datierungsgründe:}
\pagelabel{60248_dat}
Die titellose und undatierte Aufzeichnung N.\;\ref{60248} ist auf einem Träger überliefert, der weder Spuren eines Wasserzeichens noch sonstige äußere Merkmale aufweist, die eine eindeutige Datierung begründen könnten.
Das kleine Papierformat und möglicherweise auch der Schreibduktus deuten aber auf eine Abfassung des Textes nach Leibnizens Aufenthalt in Paris\indexO{Paris} 1672 bis 1676.
\pend%
\pstart%
Inhaltlich setzt sich die Aufzeichnung N.\;\ref{60248} mit Descartes'\indexN{\textls{Descartes} (Cartesius, des Cartes), René 1596\textendash1650} Erklärung vom Schwellen durchfrierenden Wassers und mit der ihr zugrundeliegenden Auffassung der fluiden Materie auseinander.
Gleich zu Beginn nimmt N.\;\ref{60248} ausführlich auf die betreffende Stelle in Descartes'\indexN{\textls{Descartes} (Cartesius, des Cartes), René 1596\textendash1650} \textit{Météores} Bezug (siehe den Erläuterungsapparat).\cite{01489}
Mit wenigstens einigen der Traktate, die Descartes\indexN{\textls{Descartes} (Cartesius, des Cartes), René 1596\textendash1650} 1637 zusammen mit dem \textit{Discours de la méthode} veröffentlicht hatte, kannte sich Leibniz schon lange.
Auf die \textit{Dioptrique}\cite{00038} verweist er bereits in der wohl im Winter 1670/1671 entstandenen \textit{Hypothesis physica nova} (Mainz 1671, §\;22; \textit{LSB} VI,\;2 N.\;40, S.\;229.2--4)\cite{00256} sowie in Aufzeichnungen über die Lichtbrechung, die sich auf das Jahr 1673 datieren lassen (vgl. etwa \textit{LSB} VIII,\;1 N.\;21\cite{02610} und N.\;22).\cite{02611}
Ob man davon ausgehen kann, dass Leibniz zu dieser Zeit auch die \textit{Météores}\cite{01489} gelesen hatte, ist unklar.
Erste explizite Erwähnungen von Descartes'\indexN{\textls{Descartes} (Cartesius, des Cartes), René 1596\textendash1650} Abhandlung, die eine Lektüre bezeugen können, finden sich in Leibnizens Brief an Hermann Conring vom 24.\;August\,/\,3.\;September 1677 (\textit{LSB} II,\;1 N.\;156, S.\;563.15\,f.)\cite{02612} sowie vornehmlich in dem Konzept eines Briefes an Unbekannte, das editorisch 1679 datiert ist (\textit{LSB} II,\;1 N.\;219, S.\;782.5\,f.).\cite{01311}
Somit kann man davon ausgehen, dass Leibniz bereits in seinen frühen Hannoveraner\indexO{Hannover} Jahren (ab 1677) die \textit{Météores}\cite{01489} hinreichend kannte, um die kritischen Bemerkungen abzufassen, die in der vorliegenden Aufzeichnung zu lesen sind.
\pend%
\pstart%
Eine Datierung der Aufzeichnung N.\;\ref{60248} auf die Zeit vor 1677 %
ist nicht auzuschließen, aber unwahrscheinlich.
Ein begründeter Terminus ante quem der Entstehung von N.\;\ref{60248} lässt sich nach heutigem Wissensstand nicht angeben.
\pend%
\end{ledgroup}%
\selectlanguage{latin}%
\frenchspacing%
\newpage%
\normalsize%
\pstart%
\noindent%
\centering%
\lbrack%
\href{\Scan{LH_37_04_093r}}%
{93\,r\textsuperscript{o}}%
\rbrack\
\hfill%
Cur aqua frigore intumescat%
\indexS{frigus}%
\indexS{aqua intumescens frigore}
\hfill%
\phantom{\lbrack93\:r\textsuperscript{o}\rbrack}%
\pend%
\vspace{0.5em}%
\pstart%
\noindent%
\edtext{%
Causam%
\indexS{causa reddita}
Cartesii\indexN{\textls{Descartes} (Cartesius, des Cartes), René 1596\textendash1650}
\textit{De meteoris} § sub finem discursus secundi%
\indexS{discursus}
reddere simulat,
sed revera non reddit.
Notat aquae particulas%
\indexS{particulae aquae}%
\indexS{aqua intumescens frigore}%
\lbrack,\rbrack\
quas concipit instar anguillarum%
\indexS{anguilla}
cum agitantur motu materiae subtilis,%
\indexS{materia subtilis}%
\indexS{motus materiae subtilis}
varie se flectere atque accomodare sibi invicem,
atque ita facilius locum replere%
\indexS{locus repletus particulis}
\lbrack;\rbrack\
cum vero
debilitato illo materiae subtilis motu%
\indexS{materia subtilis}%
\indexS{motus materiae subtilis}
non amplius ita flectuntur,
sed rigidae manent,
tunc ob curvitates earum%
\indexS{curvitas varia}
\edtext{%
varias
locum non aeque replere%
\indexS{locus repletus particulis}
ac sibi cedere posse,
ac proinde%
}{%
\lemma{varias}%
\Bfootnote{%
\textit{(1)}~sibi %
\textit{(2)}~locum non aeque replere %
\textit{(a)}~, ac %
\textit{(b)}~ac sibi cedere posse, %
\textit{(aa)}~und
\textit{(bb)}~ac proinde%
~\textit{L}%
}}
aquam frigore intumescere.%
\indexS{frigus}%
\indexS{aqua intumescens frigore}%
}{%
\lemma{Causam \lbrack...\rbrack\ intumescere}%
\Cfootnote{%
Siehe vielmehr \textsc{R.\;Descartes}, \textit{Météores}, discours I
(\textit{Discours de la méthode}, Leiden 1637, S.\;162--164;\cite{01489}
\textit{DO},\;VI, S.\;236.22--238.22).\cite{00120}%
}}
Verum hoc elusorium est.
Sane si varii arcus%
\indexS{arcus rigidus}
sive particulae curvae rigidae%
\indexS{particulae curvae rigidae}
in unum locum conjicerentur,
majus spatium occuparent %
quam si essent
\edtext{%
oblong\lbrack ae\rbrack%
}{%
\lemma{oblongi}%
\Bfootnote{%
\textit{L\;ändert Hrsg.}%
}}
et flexiles%
\indexS{particulae oblongae}%
\indexS{partculae flexiles}
instar anguillarum.%
\indexS{anguilla}
Sed cum semel partes illae flexiles
locum replevere,%
\indexS{locus repletus particulis}
superveniente induratione%
\indexS{induratio superveniens}
non magis locum majorem postu%
\lbrack%
\href{\Scan{LH_37_04_093v}}%
{93\,v\textsuperscript{o}}%
\rbrack%
labunt,
quam anguill\textlangle ae%
\indexS{anguilla}
quae\textrangle\ subito indurescere
\edtext{%
fingerentur;
situm enim quem%
}{%
\lemma{fingerentur;}%
\Bfootnote{%
\textit{(1)}~cum %
\textit{(2)}~\textlangle-- --\textrangle\ %
\textit{(3)}~situm enim quem%
~\textit{L}%
}}
proxime antehabuere retinent%
\lbrack;\rbrack\
\edtext{%
licet enim varie implicatae%
\indexS{particulae implicatae}
et arcuatae sint,%
\indexS{particulae arcuatae}%
}{%
\lemma{licet}%
\Bfootnote{%
\textit{(1)}~varie implica %
\textit{(2)}~enim varie implicatae %
\bars{et \textit{erg.}} arcuatae sint,%
~\textit{L}%
}}
tamen paulo ante
ob flexilitatem%
\indexS{flexibilitas particularum}
se sibi mutuo accommodabant.%
\indexS{particulae accomodantes sibi}
At si eaedem anguillae%
\indexS{anguilla}
jam induratae eximerentur,
ac denu\textlangle o\textrangle\ in eundem locum
\edtext{%
conjicerentur,
locus eas utique%
}{%
\lemma{conjicerentur,}%
\Bfootnote{%
\textit{(1)}~\textlangle eae\textrangle\ utique %
\textit{(2)}~locus eas utique%
~\textit{L}%
}}
capere tunc non posset.
Sed tale nihil hic contingit.
Itaque ratio Cartesii\indexN{\textls{Descartes} (Cartesius, des Cartes), René 1596\textendash1650}
sophisticata est%
\indexS{ratio sophisticata}
ex quadam festinatione profecta,%
\indexS{festinatio}
dum ille gravissimas quaestiones%
\indexS{quaestio gravis}
paucis verbis%
\indexS{verbum}
solvere posse
sibi persuadet.
\pend%
\selectlanguage{latin}%
\frenchspacing%
\normalsize%
\newpage%
